% REVISED: canonical-target reframe 2026-06-04
\section{Scope, Limits, and Price Corroboration}
\label{sec:price_scope_conclusion}

Section~\ref{sec:screening_forensics} carries the paper's central
enforcement-design result. Price enters only in a narrower,
corroborative role. It is scope evidence, speaking to whether the
award-layer screen points toward economically non-neutral
environments, and nothing more. It is \emph{not} a damages estimate,
an overcharge, a cartel markup, a causal price effect of
frequent-loser presence, or proof of a cover-bidding mechanism.
Overcharge measurement belongs to its own tradition
\citep{bryant1991price,connor2007price}, separate from screening
\citep{marshall2012economics,harrington2008detecting}. We draw the
boundary once and hold to it. The full price regressions appear in
\ref{app:scope_adaptation_submission}.

In the broad item-level regressions, frequent-loser presence is positively
associated with log negotiated unit price ($\valBetaBroad$,
$p=\valBetaBroadP$; range $\valHeadlineRange$ across estimators), and the
association runs larger in pregão than in convite. That broad sign is a
weighting artifact: holding the overlap cells fixed and moving to the
within-cell average treatment effect on the treated reverses it
($\valScopeOverlapCoef$, $p\valScopeOverlapPlt$), because the broad sample
blends a positive between-cell selection component---flagged environments sit
in structurally higher-priced cells---with a negative within-cell one. The
reversal holds outside the high-value tail, the direct-CADE overlap estimate
is null ($\valBetaCADEDirect$, $p=\valBetaCADEDirectP$), and the within-cell
movement loads on the count of genuine bidders rather than on the
frequent-loser count. A selection-driven broad imprint paired with an
unidentified within-cell channel keeps the evidence firmly on the scope side
of the line: the award-layer signal points to economically non-neutral
environments, and that is as far as it reaches. The regressions, the
decomposition table, and the bidder-count boundary are in
\ref{app:scope_adaptation_submission}.

\subsection{The Instability of a Published Administrative Cutoff}
\label{sec:adaptive_deployment_submission}

% REVISED: hostile-review action plan 2026-06-04
The FL14 cutoff is an administrative implementation rather than a structural
threshold, and Section~\ref{sec:screening_forensics} already establishes that
the ranking carries no conduct-specific residual. The warning that matters is
therefore about institutional design, not signal robustness: any
\emph{published} administrative threshold invites manipulation once firms learn
which observable patterns draw attention. Identity fragmentation and
threshold-aware capping are the most corrosive responses---simulated
manipulations of this kind can move the rank by several AUC points
(\ref{app:strategic_adaptation_submission})---and a static disclosed cutoff is
the easiest of all to game, precisely because firms can bunch against it. The
practical answer is to retire the cutoff rather than defend it. An agency
should rank firms continuously or run capacity-constrained top-$k$ queues,
refresh those queues as new participation data arrive, watch for bunching
below any operating cutoff, and follow award-layer triage with bid-layer
forensics wherever bid files can be recovered. The survivor-pool size of such a
refreshed queue is itself a point on the cost--recall frontier of
Section~\ref{sec:screening_forensics}, so investigative capacity and exposure
to gaming get chosen together rather than tuned in isolation. None of this
turns the ranking into a liability rule. It is a discipline for ordering scarce
investigative attention under the limits documented in
Sections~\ref{sec:screening_forensics} and~\ref{sec:price_scope_conclusion},
and it should never stand as proof of conduct.
